% ============================================================
% Chapter 6: Experimental Design
%   Section: Experimental Setup
% Included from chapters/06-experimental-design/experimental-design.tex
% Sources: docs/decisions/evaluation.md ("Where it runs", "Pilot");
% docs/decisions/deployment.md; the committed run.json of each report.
% ============================================================

\section{Experimental Setup}
\label{sec:setup}

\paragraph{Hardware.} Two machines are relevant. Every run reported in
this thesis was made on the development laptop: an Intel Core
i7-12700H, 16~GB of RAM and an NVIDIA RTX~3050 Laptop GPU with 4~GB, on
which about 80\% of the 3-billion-parameter model fits, the rest
running on the processor; the inference service runs on the processor.
The production target is a 4-vCPU, 16~GB virtual machine with no GPU,
on which the same model takes about 45~s per claim against about 6~s on
the laptop. Latency figures are therefore reported with the machine
they were measured on, and a CPU deployment must use the production
model limits (one call at a time, a 180~s timeout).

\paragraph{Models.} The verifier is \texttt{llama3.2:3b} (Q4\_K\_M)
served by Ollama (in production with a 4,096-token context and one
request at a time); the inference service
runs GLiNER (\texttt{gliner\_small}), the
\texttt{twitter-xlm-roberta-base-sentiment} classifier and
\texttt{BAAI/bge-m3}, with pinned library versions. Every record names
the model and the provider that answered.

\paragraph{Runs in containers.} Evaluation runs are made inside the
development stack with an extra compose layer that mounts the
evaluation data and the code (so a harness fix applies on resume
without a rebuild) and the repository read-only, so that every record
carries the commit it was made on. The layer also sets the CPU limits
the pilot needed: a 180~s model timeout, one model call at a time, a
120~s inference timeout and two claims at once, since four concurrent
claims on CPU inference pushed the ranker's embedding calls past the
30~s default.

\paragraph{Thresholds and weights.} The X-Fact runs overrode no
threshold; their effective values, written in full into each run's
manifest, were those of Section~\ref{sec:design-fact-checking}:
pertinence 0.25, ranking weights 0.45 / 0.20 / 0.20 / 0.15, pertinence
weights 0.5 / 0.5, confidence weights 0.7 / 0.3, at most five sources
per claim, one piece of evidence required for a verdict and two
independent domains for full confidence in a \texttt{TRUE}. One value
came from the environment's configuration rather than the declared
default: the duplicate threshold was 0.96 instead of 0.90. Admission is
not exercised by the harness, which checks claims, so it did not affect
these runs; it is recorded because the manifest records it. A run's
thresholds are hashed into its key.

\paragraph{Provenance.} Each report states the commits its sessions
ran on, the dataset's hash, the thresholds' hash, the harness version,
the corpus mode and point count, and the bootstrap seed. Before every
live run the health endpoint and the source probe's search check are
read and recorded, because a run's verdicts mean nothing if the search
engines refused it.

\paragraph{The split.} Before any run on the custom set, the set is to
be split once, stratified by label and language, into about one third
for tuning and two thirds for testing (Table~\ref{tab:priorities}, P0):
every threshold is tuned on the tuning part only, and every reported
number comes from the test part. \pending{the frozen split: its date,
its seed and its counts by label and language.}
